\begin{abstract}
Covert spatial cues speed and sharpen detection at their target, but a cue that points to the wrong location imposes a cost. We ask how that invalid-cue cost scales with stimulus load in a trained recurrent attention model, and whether it is a property of the top-down read-out specifically. A single convolutional, self-distilling, cross-attention recurrent agent was trained from raw pixels on a value-directed, cued orientation-change-detection task, and probed at two stimulus set sizes (four and nine locations) by reading the attention its memory pathway places on the changed location. At set size four, an invalid cue delays the attentional reallocation to the uncued change by one frame but the spotlight recovers fully thereafter (excess change-lock $+0.70$ over the uniform baseline), and behaviour confirms that uncued changes are still detected, only at higher magnitude. At set size nine the reallocation is abolished: attention never reaches the uncued change ($\approx 0$ excess at both frames), even though the cued read-out and valid-cue locking remain intact. A bottom-up element-wise-affine variant is immune, locking on uncued changes at high load but expressing a weaker cueing benefit. The set-size scaling thus traces specifically to the cue-gated top-down channel, yielding a falsifiable behavioural prediction (stated here as a prediction pending convergence) that the invalid-cue cost should grow with set size.
\end{abstract}
